Having gone through our classification of the dimension-8 operators in the previous Section we are now ready to tabulate the results.
Table~\ref{tab:summary} summarizes the results tables that follow it.
The links in the rightmost column point to the table(s) of results for a given class.
The number of types of operators and the number of Lagrangian terms in the class are given in the third and fourth columns from the left, respectively.
For comparison the number of operators from Ref.~\cite{Henning:2015alf} is given in the second column from the right.
Lines separate the classes based on the number of fermions in the class. 
Four-fermion operators are further divided into subclasses either preserving or violating baryon number.
Additionally the number to the right of the $+$ sign in $N_{\rm type, term}$ for the four-fermion operators is the number of types or terms that vanish in the absence of flavor structure
Tables~\ref{tab:smeft8class_1_2_3_4} and \ref{tab:smeft8class_5_6_7_8} contain bosonic operators.
Tables~\ref{tab:smeft8class_9}, \ref{tab:smeft8class_10_11_12_13}, \ref{tab:smeft8class_14qud}, \ref{tab:smeft8class_14le_15RR}, \ref{tab:smeft8class_15LL}, and \ref{tab:smeft8class_16_17} contain two-fermion operators.
Tables~\ref{tab:smeft8class_18_21}, \ref{tab:smeft8class_18},  \ref{tab:smeft8class_19_LL_RR}, \ref{tab:smeft8class_19_LLRR}, \ref{tab:smeft8class_19_LRRL_LRLR}, \ref{tab:smeft8class_19_slashedB}, \ref{tab:smeft8class_20_le_qu}, \ref{tab:smeft8class_20_qd_slashedB}, and \ref{tab:smeft8class_21} contain four-fermions operators.

%--------------------------------------------------------------------------------------------------------------
\begin{table}[H]
\begin{center}
\renewcommand{\arraystretch}{1.2}
\begin{tabular}[t]{c | c | c | c | c | c}
$\#$ & Class & $N_{\rm type}$ & $N_{\rm term}$ & $N_{\rm op}$~\cite{Henning:2015alf} & Table(s) \\
\hline
%--------------------------------------------------------
1    & $X^4$                & 7 & 43 & 43    & \ref{tab:smeft8class_1_2_3_4} \\
2    & $H^8$                & 1  & 1   & 1       &  \ref{tab:smeft8class_1_2_3_4} \\
3    & $H^6D^2$          & 1  & 2   & 2     &  \ref{tab:smeft8class_1_2_3_4} \\
4    & $H^4D^4$           & 1  & 3   & 3    &  \ref{tab:smeft8class_1_2_3_4} \\
5    & $X^3H^2$          & 3  & 6   & 6       & \ref{tab:smeft8class_5_6_7_8} \\
6    & $X^2H^4$          & 5 & 10 & 10    & \ref{tab:smeft8class_5_6_7_8} \\
7    & $X^2H^2D^2$    & 4 & 18 & 18   & \ref{tab:smeft8class_5_6_7_8} \\
8    & $XH^4D^2$        &  2 & 6   & 6     & \ref{tab:smeft8class_5_6_7_8} \\ \hline
%--------------------------------------------------------
9    & $\psi^2X^2H$     & 16 & 96 & $96 n_g^2$ & \ref{tab:smeft8class_9} \\
10  & $\psi^2XH^3$     & 8 & 22 & $22 n_g^2$  &  \ref{tab:smeft8class_10_11_12_13} \\
11  & $\psi^2H^2D^3$  & 6 & 16 & $16 n_g^2$ &  \ref{tab:smeft8class_10_11_12_13} \\
12  & $\psi^2H^5$       & 3   & 6   & $6 n_g^2$    &  \ref{tab:smeft8class_10_11_12_13} \\
13  & $\psi^2H^4D$    & 6 & 13 & $13 n_g^2$  &  \ref{tab:smeft8class_10_11_12_13} \\
14  & $\psi^2X^2D$     & 21  & 57 & $57 n_g^2$  &  \ref{tab:smeft8class_14qud}, \ref{tab:smeft8class_14le_15RR} \\
15  & $\psi^2XH^2D$  & 16 & 92  & $92 n_g^2$ &  \ref{tab:smeft8class_14le_15RR}, \ref{tab:smeft8class_15LL} \\
16  & $\psi^2XHD^2$  & 8  & 48  & $48 n_g^2$ &  \ref{tab:smeft8class_16_17} \\
17  & $\psi^2H^3D^2$  & 3 & 36  & $36 n_g^2$ &  \ref{tab:smeft8class_16_17} \\ \hline 
%--------------------------------------------------------
18$(B)$  & \multirow{2}{*}{$\psi^4H^2$}  & 19  & 75 + 1 & $n_g^2 (67 n_g^2 + n_g +7)$  &  \ref{tab:smeft8class_18_21}, \ref{tab:smeft8class_18} \\
18$(\slashed B)$  &                                 & $4+3$  & $12+8$ & $\tfrac{1}{3} n_g^2 (43 n_g^2 - 9 n_g + 2)$  &  \ref{tab:smeft8class_18_21} \\ 
19$(B)$  & \multirow{2}{*}{$\psi^4X$}  & $40+5$  & $156+12$ & $4 n_g^2 (40 n_g^2 - 1)$ &  \ref{tab:smeft8class_19_LL_RR}, \ref{tab:smeft8class_19_LLRR}, \ref{tab:smeft8class_19_LRRL_LRLR} \\
19$(\slashed B)$  &                             & 4  & $44+12$   & $2 n_g^3 (21 n_g + 1)$    &  \ref{tab:smeft8class_19_slashedB} \\
20$(B)$  & \multirow{2}{*}{$\psi^4HD$}  & 16  & $134+2$ & $n_g^3 (135 n_g - 1)$ &  \ref{tab:smeft8class_20_le_qu}, \ref{tab:smeft8class_20_qd_slashedB} \\
20$(\slashed B)$  &                                 & 7 & 32 & $n_g^3 (29 n_g + 3)$    &  \ref{tab:smeft8class_20_qd_slashedB} \\ 
21$(B)$  & \multirow{2}{*}{$\psi^4D^2$}  & 18  & 55 & $\tfrac{11}{2} n_g^2 (9 n_g^2 + 1)$ &  \ref{tab:smeft8class_18_21}, \ref{tab:smeft8class_21} \\
21$(\slashed B)$  &                                 & 4  & $10+2$ & $n_g^3 (11 n_g - 1)$   &  \ref{tab:smeft8class_18_21} \\ \hline \hline
%--------------------------------------------------------
  & $B$                & $204+5$  & $895+15$ & 895(36971), $n_g = 1 (3)$ & \\
  & $\slashed B$  & $19+3$     & $98+22$   & 98(7836), $n_g = 1 (3)$   & \\
  & Total               & $223+8$  & $993+37$ & 993(44807), $n_g = 1 (3)$ &
\end{tabular}
\end{center}
\caption{Summary of the contents of the tables to follow. 
The links in the rightmost column point to the table(s) of results for a given class.
The number of types of operators and the number of Lagrangian terms in the class are given in the third and fourth columns from the left, respectively.
For comparison the number of operators from Ref.~\cite{Henning:2015alf} is given in the second column from the right.
Lines separate the classes based on the number of fermions in the class. 
Four-fermion operators are further divided into subclasses either preserving or violating baryon number.
Additionally the number to the right of the $+$ sign in $N_{\rm type(term)}$ for the four-fermion operators is the number of types(terms) that vanish in the absence of flavor structure.}
\label{tab:summary}
\end{table}